\subsection[chaîne de numérisation]{La chaîne de numérisation de CujasNum}

La numérisation des documents anciens demande au préalable une réflexion poussée \enquote{portant sur les choix techniques comme intellectuels}. Pour la bibliothèque Cujas, l'opération de numérisation représente en effet un puissant levier pour valoriser scientifiquement ses collections patrimoniales et en élargir durablement l'accès.

Malgré la richesse des collections conservées dans la réserve de la bibliothèque Cujas, celle-ci ne rencontre pas immédiatement le succès escompté auprès des usagers. Frédérique Joannic-Seta corrobore ce constat, attribuant ce manque d'attrait à \enquote{ses difficultés d'ouverte de consultation qui entraînent une faible utilisation au regard de sa réelle richesse}. Par ailleurs, de nombreux volumes se trouvent alors en mauvais état, ayant subi aussi bien des \enquote{inondations exceptionnelles dans les magasins que des dégradations quotidiennes suscitées par les manipulations et les conditions de cnoservation, longtemps médiocres}. 

En 1998, soucieuse de valoriser cette réserve patrimoniale, la bibliothèque Cujas décide d'engager la numérisation de ses collections. Son choix est dicté par \enquote{un souci de conservation des documents par la mise à disposition des lecteurs d'un support de substitution de qualité}. Le second objectif vise à faciliter l'accès aux fonds pour permettre aux lecteurs \enquote{une meilleure appropriation du texte en leur proposant des facilités de reproduction}. 

A la suite de l'exposition virtuelle consacrée à Jean Carbonnier, la bibliothèque Cujas dispose désormais d'une chaîne de numérisation en interne opérationnelle et prête à être mobilisée. Cette infrastructure s'est d'ailleurs renforcée grâce à la création de la célèbre \enquote{liste Pfister-Roumy}, une campagne de numérisation \enquote{alimentée enti-èrement en interne grâce aux compétences du département informatique}. Pour ce corpus, l'établissement fait le choix de traiter en interne exclusivement les \enquote{ouvrages postérieurs à 1830 avec l'océrisation et la relecture partielle sur les parties structurantes}. A l'inverse, pour les documents antérieurs à 1830, la bibliothèque Cujas privilégie une externatisation de la numérisation auprès d'un partenaire spécialisé, la société Puces et Plumes. Cette décise vise à \enquote{préserver le document}, car le scanner représente un risque pour ses ouvrages déjà fragiles. Le prestataire se charge de numériser les documents en \enquote{mode image, de faire la saisie fines des tables des matières}. 

Afin de garantir, la bonne exécution de la campagne de numérisation, la méthodologie de travail de la bibliothèque Cujas s'appuie sur l'établissement systématique d'un bordereau constituer de différentes parties. La première partie correspond à la notice bibliographique du document; la seconde détaille précisement sa pagination ainsi que les éventuelles particularités observées, par exemple des pages portant la même pagination; enfin, la dernière partie dresse l'état du document avant le transfert chez le prestataire.

Une fois la numérisation effectuée, la bibliothèque réceptionnne les images réalisées sous le format TIFF et effectue \enquote{un contrôle physique, livre en main, pour vérifier la complétude des images et s'assurer qu'aucune page ne manque par rapport au bordereau}. C'est par la suite qu'un script informatique permet de générer les formats de diffusion (JPEG, PDF) et de produire la table des correspondance \enquote{qui fait correspondre chaque numéro d'image à la pagination réelle du livre}. Simultanément, tandis que la personne en charge du contrôle qualité vérifie la conformité des informations renseignées dans le Dublin Core, des mentors étudiants se voient confier la saisie intégrale des tables des matières, une tâche essentielle pour restituer fidèlement la structure intellectuelle de l'ouvrage. 

La validation de ce contrôle qualité ouvre la voie à la phase finale de publication, qui requiert de minutieuses vérifications formelles et sémantiques. Cette étape consiste à vérifier \enquote{la résolution et la compression des images JPEG et TIFF dans IrfanView, de vérouiller les PDF grâce à la protection par un mot de passe}, et à valider la conformité de la structure XML du Dublin Core afin de conjurer tout risque de rejet par le CINES. Une fois ces critères validés, les données sont alors \enquote{transférées via un cleint SFTP sécurisé durant la nuit vers l'espace de publication}. Ce protocole de transfert différé et restreint a été mis en oeuvre à la suite d'une tentative d'attaque par rançonnage subie par l'établissement en 2017. Pour sécuriser les données, le département information technique (DIT) a donc mis en place l'utilisation de ce protocole. 

Au démarrage du traitement numérique, la Reconnaissance Optique de Caractères (OCR) est appliquées exclusivement aux documents postérieurs à 1830 afin de permettre la recherche plein-texte. Pour les ouvrages antérieurs à cette date, la numérisation s'effectue uniquement en mode image pour préserver les supports physiques. l'absence d'OCR sur le corpus de ces textes anciens est palliée par l'ajout sémantique de la table des matières, garantissant ainsi une navigation intellectuelle dans le document.

Une autre particularité liée à la chaîne de numérisation réside dans l'indexation appliquée aux documents numériques de la bibliothèque numérique. C'est l'indexation Hauville qui est crée à partir de l'indexation RAMEAU. 